\chapter{Estudo de viabilidade técnico-econômica}



Com base nas análises realizadas, foram identificadas oportunidades significativas para a melhoria da eficiência no uso da energia elétrica nas instalações da unidade consumidora, especialmente nos sistemas de \textbf{iluminação, climatização, acionamentos motorizados} e \textbf{geração própria de energia}. A implementação das ações propostas poderá resultar em uma economia estimada de \textbf{\textcolor{red}{[\_\_\_\_]}}~MWh/ano, representando uma redução de aproximadamente \textbf{\textcolor{red}{[\_\_\_\_]}}~\% no consumo total de energia elétrica da unidade.

Além disso, estima-se uma redução da demanda contratada da ordem de \textbf{\textcolor{red}{[\_\_\_\_]}}~kW, correspondente a cerca de \textbf{\textcolor{red}{[\_\_\_\_]}}~\% da demanda máxima registrada no período analisado.

Para viabilizar a execução do projeto, foi previsto um investimento total estimado em \textbf{R\$~\textcolor{red}{[\_\_\_\_]}}. Esse valor contempla a aquisição de equipamentos eficientes, serviços de instalação e comissionamento, mão de obra especializada, bem como os encargos indiretos associados (BDI – Benefícios e Despesas Indiretas).

A análise de viabilidade econômico-financeira foi conduzida com base na construção de um fluxo de caixa projetado para um horizonte de \textbf{5 anos}, utilizando os seguintes indicadores de desempenho:

\begin{itemize}
    \item \textbf{Valor Presente Líquido (VPL):} R\$~\textcolor{red}{\textbf{[\_\_\_\_]}}
    \item \textbf{Taxa Interna de Retorno (TIR):} \textcolor{red}{\textbf{[\_\_\_\_]}}~\%
    \item \textbf{Tempo de Retorno Simples (Payback):} \textcolor{red}{\textbf{[\_\_\_\_]}}~anos
    \item \textbf{Relação Custo-Benefício (RCB):} \textcolor{red}{\textbf{[\_\_\_\_]}}
\end{itemize}

Esses indicadores foram obtidos a partir da estimativa de economia financeira anual gerada pelas medidas propostas, levando em consideração o custo médio atual da energia, a depreciação dos equipamentos e a ausência de receitas com venda de excedentes (no caso da geração fotovoltaica compensada via autoconsumo).

De acordo com os critérios estabelecidos no \textbf{Módulo 7 do Programa de Eficiência Energética da ANEEL (PROPEE)}~\cite{propee_aneel}, que orienta a estruturação e avaliação dos projetos com base na relação custo-benefício e horizonte de retorno, os resultados obtidos indicam que o projeto apresenta \textbf{viabilidade técnico-econômica favorável} para implantação, atendendo aos requisitos normativos e aos parâmetros usuais de análise do setor.

